% =====================================================================
\chapter{Problemas}
\label{cap:problemas}
% =====================================================================

Cinco problemas de 12 puntos cada uno (60 puntos en total), uno por módulo. La puntuación de cada apartado figura entre paréntesis. Los problemas marcados con \emph{Fuente} están adaptados de exámenes, hojas de problemas o trabajos de curso de universidades; el enunciado es de redacción propia y los datos se han cambiado para que puedan resolverse con calculadora. Tres de ellos (B.1, B.3 y B.5) proceden de exámenes o de colecciones de preparación de exámenes. Las soluciones completas están en el \cref{ap:soluciones}.

% ---------------------------------------------------------------------
\section*{Problema B.1. Hornos, costes fijos y una disyunción (M24)}
\phantomsection\addcontentsline{toc}{section}{Problema B.1. Hornos, costes fijos y una disyunción}

\begin{ejercicio}{2}{eje:milp}
Una planta de tratamiento térmico debe procesar exactamente $D=100$~t en un turno con tres hornos. Encender el horno $i$ cuesta $f_i$ (precalentamiento y energía fija); procesar cada tonelada cuesta $c_i$ en energía; si se enciende, por estabilidad del proceso debe tratar al menos $L_i$ toneladas, y nunca más de su capacidad $U_i$:
\begin{center}\small
\begin{tabular}{@{}lcccc@{}}
\toprule
Horno & $f_i$ (€) & $c_i$ (€/t) & $L_i$ (t) & $U_i$ (t) \\
\midrule
1 & 400 & 12 & 30 & 80 \\
2 & 150 & 16 & 20 & 50 \\
3 & 250 & 14 & 25 & 60 \\
\bottomrule
\end{tabular}
\end{center}
\begin{enumerate}[label=(\alph*)]
  \item \pt{3} Formule el problema de coste mínimo como un programa lineal entero mixto con variables $x_i\ge 0$ (toneladas) e $y_i\in\{0,1\}$ (encendido), usando restricciones de tipo big-M. Indique qué representa cada restricción.
  \item \pt{2} Determine la menor constante $M_i$ válida de cada restricción $x_i\le M_iy_i$ y justifique que no elimina ninguna solución factible.
  \item \pt{3} Resuelva el problema enumerando los patrones de encendido $y\in\{0,1\}^{3}$; para cada patrón factible, resuelva a mano el programa lineal en $x$. Dé el plan óptimo y su coste.
  \item \pt{2} Calcule el valor de la relajación lineal ($y_i\in[0,1]$) con las $M_i$ de~(b) y el \emph{gap} de integralidad. Repita el cálculo si se escribe la capacidad como restricción aparte, $x_i\le U_i$, y el enlace con $M=10^{4}$ para los tres hornos. ¿Qué concluye?
  \item \pt{2} La acometida eléctrica impone que, si el horno~2 está encendido, el horno~1 no trate más de 60~t. Escriba la restricción lineal con la menor constante válida y obtenga el nuevo óptimo.
\end{enumerate}
\fuente{apartados~(a), (b) y~(e) adaptados de MIT 15.053, primavera de 2013, colección de problemas de preparación del segundo examen parcial (\emph{Midterm Exam~2: Practice Problem Set}), problema~2 (apertura de cajas de cobro con coste fijo y restricciones de enlace) y problema~4(a) (restricción disyuntiva «al menos una de dos»)~\cite{mit15053mid}; apartado~(d) inspirado en la hoja de problemas~5, problema~4, del mismo curso (elegir la menor $M$ válida)~\cite{mit15053ps}.}
\end{ejercicio}

% ---------------------------------------------------------------------
\section*{Problema B.2. Un paso de optimización bayesiana a mano (M25)}
\phantomsection\addcontentsline{toc}{section}{Problema B.2. Un paso de optimización bayesiana a mano}

\begin{ejercicio}{2}{eje:bo}
Se ajusta la temperatura de curado de un adhesivo (variable codificada $x$) para minimizar el tiempo de ciclo tipificado. Se modeliza $f\sim\mathcal{GP}(0,k)$ con núcleo exponencial cuadrático $k(x,x')=\exp\bigl(-\tfrac12(x-x')^{2}/\ell^{2}\bigr)$, $\sigma_f=1$ y $\ell=1$. Se han hecho dos ensayos sin ruido: $f(0)=0{,}5$ y $f(2)=-0{,}8$.
\begin{enumerate}[label=(\alph*)]
  \item \pt{3} Escriba $K$, calcule $K^{-1}$ y $\alpha=K^{-1}y$, y obtenga la media y la varianza posteriores en $x=1$.
  \item \pt{2} Obtenga la media y la varianza posteriores en $x=3$.
  \item \pt{3} Calcule la probabilidad de mejora y la mejora esperada en $x=1$ y en $x=3$, y la cota inferior de confianza $\mu-2\sigma$ en ambos. ¿Qué punto se ensaya a continuación con cada criterio? Explique el resultado en términos de explotación y exploración.
  \item \pt{2} Si los ensayos tuvieran ruido gaussiano de varianza $\sigma_\varepsilon^{2}=0{,}1$, calcule la varianza posterior de $f(2)$ y la de una nueva observación $y(2)$. ¿Cuál debe usarse en la mejora esperada y por qué?
  \item \pt{2} Sin ruido, se cambia la longitud de escala a $\ell=0{,}3$. Calcule aproximadamente $\mu(1)$, $\sigma^{2}(1)$ y $\EI(1)$, y explique cómo afecta la longitud de escala a la decisión.
\end{enumerate}
\fuente{adaptado del trabajo de curso~1, «Gaussian processes», de Cambridge 4F13 \emph{Probabilistic Machine Learning} (Michaelmas de 2024), que pide predecir con un proceso gaussiano, interpretar la incertidumbre y estudiar el efecto de los hiperparámetros~\cite{cam4f13}, y del ejemplo de la mejora esperada de Jones, Schonlau y Welch~\cite{jones}.}
\end{ejercicio}

% ---------------------------------------------------------------------
\section*{Problema B.3. Riccati, MPC de un paso y Bellman con coste lineal (M26)}
\phantomsection\addcontentsline{toc}{section}{Problema B.3. Riccati, MPC de un paso y Bellman con coste lineal}

\begin{ejercicio}{3}{eje:mpc}
La desviación de temperatura de un depósito de fluido térmico respecto a su consigna, $x_k$ (K), evoluciona con $x_{k+1}=a\,x_k+b\,u_k$, donde $u_k$ es la potencia de calefacción o refrigeración (kW, en variables de desviación).
\begin{enumerate}[label=(\alph*)]
  \item \pt{3} Con $a=0{,}9$, $b=0{,}5$ K/kW, coste $\sum_{j=0}^{1}(Qx_j^{2}+Ru_j^{2})+P_fx_2^{2}$, $Q=1$, $R=0{,}5$, $P_f=1$ y $x_0=2$~K, aplique la recursión de Riccati para obtener $K_1$, $P_1$, $K_0$ y $P_0$, la secuencia óptima sin restricciones y su coste. Compruebe que coincide con $P_0x_0^{2}$.
  \item \pt{3} El actuador está limitado a $\abs{u_j}\le 1$~kW. Escriba el coste como $J(U)=U\T HU+2x_0F\T U+\text{cte}$, resuelva el MPC de horizonte 2, indique qué decisión se aplica en el instante actual y calcule el multiplicador de la restricción activa.
  \item \pt{4} En otra operación se usa el coste descontado con término lineal
  \[
  J=\sum_{t=0}^{\infty}\gamma^{t}\bigl(q\,x_t+r\,u_t^{2}\bigr),\qquad q=2,\ r=1,\ \gamma=0{,}95,
  \]
  sin restricciones y con los mismos $a$ y $b$. Plantee la ecuación de Bellman para $V(z)=\min J$ con $x_0=z$, compruebe que una función afín $V(z)=pz+s$ la satisface y calcule $p$, $s$ y la acción óptima.
  \item \pt{2} Se añade a~(c) la restricción $\abs{u_t}\le 2$~kW. Explique por qué $V$ sigue siendo afín con la misma pendiente $p$, y calcule la nueva acción óptima y el nuevo término $s$.
\end{enumerate}
\fuente{apartado~(c) adaptado de Stanford EE363 \emph{Linear Dynamical Systems} (S.~Boyd), examen final de invierno de 2006, problema~8, «LQR with linear state cost», escalarizado~\cite{ee363final}; apartados~(a) y~(b) siguen la clase~16 de EE363 y el capítulo~1 de Rawlings, Mayne y Diehl~\cite{ee363,rmd}.}
\end{ejercicio}

% ---------------------------------------------------------------------
\section*{Problema B.4. Sensibilidad adjunta y SIMP en una barra (M27)}
\phantomsection\addcontentsline{toc}{section}{Problema B.4. Sensibilidad adjunta y SIMP en una barra}

\begin{ejercicio}{3}{eje:simp}
Una barra empotrada en su extremo izquierdo (nodo~0) se discretiza en tres elementos en serie con nodos libres~1, 2 y~3. El elemento $e$ tiene rigidez axial $k_e=x_e^{\,p}k_0$, con $k_0=1000$~N/mm, $p=3$ y densidades $x=(1{,}0;\ 0{,}8;\ 0{,}5)$ ($\Emin$ despreciable). Actúan cargas de 100~N en los nodos~2 y~3, es decir, $F=(0,100,100)$~N.
\begin{enumerate}[label=(\alph*)]
  \item \pt{2} Ensamble la matriz de rigidez $K\in\R^{3\times 3}$ (tridiagonal), obtenga los desplazamientos $U$ y la flexibilidad $c=F\T U$.
  \item \pt{3} Calcule $\partial c/\partial x_e$ para los tres elementos con el método adjunto, justificando que el problema es autoadjunto. Compruebe uno de los valores derivando directamente la expresión de $c$.
  \item \pt{3} Un sensor mide el desplazamiento del nodo~2 y se quiere controlar $J=u_2=e_2\T U$. Plantee y resuelva el problema adjunto y calcule $\partial J/\partial x_e$. Interprete el valor obtenido para el elemento~3.
  \item \pt{2} Se aplica una actualización por criterios de optimalidad con $\eta=1/2$, límite de movimiento $m=0{,}2$ y volumen fijo $\sum_ex_e=2{,}3$. Calcule $x^{\text{nuevo}}$ con $\Lambda=300$, compruebe si cumple el volumen y diga hacia dónde debe moverse $\Lambda$. Verifique que $\Lambda=270$ cumple la restricción.
  \item \pt{2} Con $p=1$ el problema es convexo. Halle con las condiciones de Karush--Kuhn--Tucker las densidades que minimizan $c$ con $\sum_ex_e=2{,}3$ (ignore las cotas $x_e\le 1$ y compruebe después que se cumplen) y compare la flexibilidad con la del diseño inicial evaluado con $p=1$.
\end{enumerate}
\fuente{apartados~(b) y~(c) adaptados de MIT 18.S096 \emph{Matrix Calculus for Machine Learning and Beyond} (IAP 2023), hoja de problemas~2, problema~1 (gradiente de una función de la solución de un sistema tridiagonal con una sola resolución adjunta)~\cite{mit18s096ps2}; apartados~(d) y~(e) siguen el capítulo~1 de Bendsøe y Sigmund~\cite{bs}.}
\end{ejercicio}

% ---------------------------------------------------------------------
\section*{Problema B.5. ¿Compensa el optimizador? Intervalos y valor actual (M28)}
\phantomsection\addcontentsline{toc}{section}{Problema B.5. ¿Compensa el optimizador? Intervalos y valor actual}

\begin{ejercicio}{2}{eje:valor}
Una planta ha probado durante $n=16$ semanas un optimizador de planificación energética. Cada semana se calcula, con el gemelo digital calibrado y los mismos pedidos, precios y averías, la diferencia de coste $d_i=$ (coste con el plan actual) $-$ (coste con el optimizador). Se obtiene $\bar d=1200$~€/semana y $s_d=1600$~€/semana.
\begin{enumerate}[label=(\alph*)]
  \item \pt{3} Calcule el intervalo de confianza del 95\,\% para la mejora media semanal y expréselo en €/año (52 semanas). Enuncie las hipótesis que lo justifican.
  \item \pt{2} Por separado, los costes semanales de cada política tienen desviación típica de 3000~€ y correlación $0{,}86$. Calcule la desviación típica de la diferencia pareada y la que tendría la diferencia con semanas independientes. ¿Cuántas veces más semanas harían falta sin emparejar para la misma precisión?
  \item \pt{2} Con $B=1000$ remuestras \emph{bootstrap} de las 16 diferencias, los percentiles 5 y 95 de las medias remuestreadas son 560 y 1890~€/semana. Dé los intervalos \emph{bootstrap} percentil y básico del 90\,\%. ¿Por qué pueden no ser simétricos respecto a $\bar d$?
  \item \pt{3} La inversión es $I=60\,000$~€ y la operación cuesta $12\,000$~€/año; el horizonte es de 5 años con $\rho=8\,\%$. Calcule el VAN con la mejora media y con el extremo inferior del intervalo de~(a), y la mejora semanal mínima que hace $\VAN=0$. ¿Puede afirmarse con un 95\,\% de confianza que el proyecto es rentable?
  \item \pt{2} ¿Cuántas semanas pareadas harían falta para detectar, con $\alpha=0{,}05$ bilateral y potencia $0{,}8$, una mejora media igual a la de rentabilidad nula de~(d), suponiendo $\sigma_D=1600$~€?
\end{enumerate}
\fuente{adaptado de MIT 18.05 \emph{Introduction to Probability and Statistics} (Orloff y Bloom), examen final de primavera de 2022, problemas~9 (contraste $t$ e intervalo de confianza), 11 (potencia de un contraste) y~13 (intervalo \emph{bootstrap} percentil)~\cite{mit1805final}; apartado~(d), de la sesión «Present Value Relations» de MIT 15.401~\cite{mit15401}.}
\end{ejercicio}
